\section{CGGMP21 for ECDSA}

\label{sec:cggmp}
% ============================================================================

ECDSA threshold signing is more complex than Schnorr because the ECDSA
equation $s = k^{-1}(H(m) + r \cdot x)$ involves multiplication of secret
values ($k^{-1}$ and $x$), requiring MPC multiplication protocols.

\subsection{Key Generation}

CGGMP21~\cite{cggmp21} key generation extends Feldman VSS with Paillier
encryption for MPC multiplication:
\begin{enumerate}[leftmargin=1.4em]
    \item Each party generates a Paillier key pair $(N_i, \phi_i)$ with
    $N_i > 2^{2048}$ and proves correctness via a zero-knowledge proof of
    safe prime product.
    \item Shamir sharing of the ECDSA private key $x$ proceeds as in FROST DKG.
    \item Each party encrypts their share under their Paillier key and
    broadcasts the ciphertext with a range proof.
\end{enumerate}

\subsection{Threshold Signing Protocol}

The CGGMP21 signing protocol consists of four phases:

\begin{enumerate}[leftmargin=1.4em]
    \item \textbf{Nonce generation}: Each signer $P_i$ samples $k_i, \gamma_i$
    and broadcasts commitments $\Gamma_i = \gamma_i \cdot G$ and Paillier
    encryptions $K_i = \texttt{Enc}(k_i)$.
    \item \textbf{MPC multiplication}: Parties compute $\delta_i = k_i \gamma_j
    + k_j \gamma_i$ pairwise using the Paillier homomorphism, with
    affine-group range proofs ensuring honest behavior.
    \item \textbf{Nonce reveal}: Parties reveal $\delta_i$ and reconstruct
    $\delta = k \cdot \gamma$. Compute $R = (\delta^{-1}) \cdot \Gamma$ where
    $\Gamma = \sum_i \Gamma_i$. Extract $r = R.x \bmod q$.
    \item \textbf{Partial signatures}: Each party computes $\sigma_i = k_i
    \cdot H(m) + k_i \cdot r \cdot \chi_i$ where $\chi_i = x_i \cdot
    \lambda_i$ is the Lagrange-weighted share.
\end{enumerate}

\begin{theorem}[Identifiable Abort]
If any party deviates from the CGGMP21 protocol, the honest parties can
identify the cheater with probability 1 and exclude them from future signing
sessions. This relies on the zero-knowledge range proofs attached to every
Paillier operation.
\end{theorem}

% ============================================================================
